\documentclass[12pt,a4paper]{article}

% ---------------------------------------------------------------
% PACKAGES
% ---------------------------------------------------------------
\usepackage{amsmath}
\usepackage{amssymb}
\usepackage{geometry}
\geometry{margin=2.5cm}
\usepackage{booktabs}
\usepackage{xcolor}
\usepackage{hyperref}
\usepackage{listings}
\usepackage{pgfplots}
\pgfplotsset{compat=1.18}
\usepackage{tcolorbox}
\tcbuselibrary{skins, breakable}
\usepackage{enumitem}
\usepackage{fancyhdr}
\usepackage{tikz}
\usepackage{array}

% ---------------------------------------------------------------
% TCOLORBOX ENVIRONMENTS
% ---------------------------------------------------------------
\newtcolorbox{curiositybox}[1]{colback=yellow!10, colframe=orange!80!black, fonttitle=\bfseries, title=#1, breakable}
\newtcolorbox{infobox}[1]{colback=blue!5, colframe=blue!60!black, fonttitle=\bfseries, title=#1, breakable}
\newtcolorbox{mistakebox}[1]{colback=red!5, colframe=red!60!black, fonttitle=\bfseries, title=#1, breakable}
\newtcolorbox{learnbox}[1]{colback=green!5, colframe=green!60!black, fonttitle=\bfseries, title=#1, breakable}

% ---------------------------------------------------------------
% LSTLISTING SETUP
% ---------------------------------------------------------------
\lstset{language=Python, basicstyle=\ttfamily\small, keywordstyle=\color{blue},
        commentstyle=\color{gray}, stringstyle=\color{red},
        showstringspaces=false, breaklines=true, frame=single}

% ---------------------------------------------------------------
% FANCYHDR
% ---------------------------------------------------------------
\pagestyle{fancy}
\fancyhf{}
\lhead{Topic 17: Residues}
\rhead{Unit II --- Complex Variables}
\cfoot{\thepage}

% ---------------------------------------------------------------
% TITLE
% ---------------------------------------------------------------
\title{\textbf{Topic 17: Residues}\\
\large Unit II --- Complex Variable: Integration\\
Complex Variables and Special Functions\\
JNTUA College of Engineering (Autonomous)}
\author{}
\date{}

% ---------------------------------------------------------------
\begin{document}
\maketitle
\tableofcontents
\newpage

% ================================================================
% OPENING HOOK
% ================================================================
\begin{curiositybox}{Hook}
A Laurent series may contain infinitely many terms, but when contour integration enters the scene,
one coefficient suddenly becomes the star of the show --- the coefficient of $\dfrac{1}{z - z_0}$.
That coefficient is called the \textbf{residue}, and it can determine the entire contour integral.
\end{curiositybox}

% ================================================================
\section{Definition of Residue}
% ================================================================

\begin{infobox}{Definition: Residue at a Singular Point}
Let $f(z)$ be analytic in a punctured neighbourhood $0 < |z - z_0| < R$ of a singular point $z_0$.
The Laurent expansion of $f$ about $z_0$ is
\[
   f(z) = \sum_{n=-\infty}^{\infty} a_n (z - z_0)^n,
\]
where the coefficients are given by
\[
   a_n = \frac{1}{2\pi i} \oint_{C} \frac{f(z)}{(z-z_0)^{n+1}}\,dz.
\]
The \textbf{residue} of $f$ at $z_0$ is the coefficient $a_{-1}$:
\[
   \operatorname{Res}(f, z_0) = a_{-1}.
\]
Equivalently,
\[
   \operatorname{Res}(f, z_0)
   = \frac{1}{2\pi i}\oint_{C} f(z)\,dz,
\]
where $C$ is any simple closed positively oriented contour encircling $z_0$ and no other singularity.
\end{infobox}

The residue is a single complex number attached to a singularity. Its importance lies in the fact that
the entire closed-contour integral of $f$ around an isolated singularity collapses to $2\pi i$ times
this one number, regardless of how complicated the rest of the Laurent series may be.

\begin{learnbox}{Key Idea --- Section 1}
$\operatorname{Res}(f,z_0) = a_{-1}$ is the coefficient of $(z-z_0)^{-1}$ in the Laurent
series of $f$ about $z_0$.  All higher-order Laurent terms integrate to zero around a closed
contour; only the $a_{-1}$ term survives.
\end{learnbox}

% ================================================================
\section{Why Residues Matter}
% ================================================================

Before residues, evaluating $\oint_C f(z)\,dz$ required parameterising the contour and
performing a direct integration --- a process that can be extremely tedious.
Residues reduce this to \emph{algebraic arithmetic}: identify isolated singularities inside $C$,
compute a number at each, sum, and multiply by $2\pi i$.

\begin{center}
\begin{tabular}{@{}lll@{}}
\toprule
\textbf{Situation} & \textbf{Without Residues} & \textbf{With Residues} \\
\midrule
$\oint_C \frac{e^z}{z^3}\,dz$ & Parameterise $C$, expand $e^z$, integrate term-by-term & Read off $a_{-1} = \tfrac{1}{2}$; answer $= \pi i$ \\
$\oint_C \frac{1}{z^2+1}\,dz$ & Partial fractions then parameterise & $\operatorname{Res} = \tfrac{1}{2i}$; answer $= \pi$ \\
Real integral $\int_0^{2\pi}\!\frac{d\theta}{2+\cos\theta}$ & No elementary method & Convert to contour, read residue \\
\bottomrule
\end{tabular}
\end{center}

\begin{learnbox}{Key Idea --- Section 2}
Residues turn closed-contour complex integrals into pure algebra.  They also unlock otherwise
intractable real definite integrals via the substitution $z = e^{i\theta}$ or by closing contours
in the upper/lower half-plane.
\end{learnbox}

% ================================================================
\section{Residue Formulas}
% ================================================================

\begin{infobox}{Formula 1 --- Coefficient from Laurent Series}
Expand $f(z)$ in a Laurent series about $z_0$ and read off the $a_{-1}$ coefficient directly.
\[
   \operatorname{Res}(f, z_0) = a_{-1}
   \qquad \text{(from the Laurent expansion)}.
\]
This method is always correct but can be laborious when the expansion is not obvious.
\end{infobox}

\begin{infobox}{Formula 2 --- Simple Pole ($m=1$)}
If $z_0$ is a simple pole of $f$, then
\[
   \operatorname{Res}(f, z_0) = \lim_{z \to z_0}(z - z_0)\,f(z).
\]
\end{infobox}

\begin{infobox}{Formula 3 --- Pole of Order $m$}
If $z_0$ is a pole of order $m \geq 1$, then
\[
   \operatorname{Res}(f, z_0)
   = \frac{1}{(m-1)!}\lim_{z \to z_0}
     \frac{d^{\,m-1}}{dz^{\,m-1}}\!\left[(z-z_0)^m f(z)\right].
\]
\end{infobox}

\begin{infobox}{Formula 4 --- Rational Function $f = p(z)/q(z)$ at a Simple Zero of $q$}
If $q(z_0) = 0$, $q'(z_0) \neq 0$, and $p(z_0) \neq 0$, then $z_0$ is a simple pole of $p/q$ and
\[
   \operatorname{Res}\!\left(\frac{p}{q},\, z_0\right) = \frac{p(z_0)}{q'(z_0)}.
\]
This is the most efficient formula for rational functions with simple poles.
\end{infobox}

\begin{learnbox}{Key Idea --- Section 3}
Four formulas exist.  The Laurent-series method is universal.  For simple poles of rational
functions, the $p/q'$ formula is fastest.  For higher-order poles, the $(m-1)$th derivative formula
is the standard route.
\end{learnbox}

% ================================================================
\section{Choosing the Right Method}
% ================================================================

Selecting the wrong formula wastes time or produces errors.  The following decision guide
covers every common case:

\begin{enumerate}[leftmargin=*]
  \item \textbf{Identify the type and order of the singularity first.}
        Factor the denominator or examine the Laurent series behaviour near $z_0$.
  \item \textbf{Simple pole ($m=1$):}
        Use $\lim_{z\to z_0}(z-z_0)f(z)$.
        If $f = p/q$ and $q'(z_0)\neq 0$, prefer $p(z_0)/q'(z_0)$ --- one less limit to evaluate.
  \item \textbf{Pole of order $m \geq 2$:}
        Multiply by $(z-z_0)^m$, differentiate $(m-1)$ times, divide by $(m-1)!$, then take $z\to z_0$.
  \item \textbf{Essential singularity or when $m$ is unclear:}
        Expand $f$ in a Laurent series directly and read $a_{-1}$.
  \item \textbf{Trigonometric or exponential numerators:}
        First simplify using Taylor expansions ($\sin z = z - z^3/6 + \cdots$, $e^z = 1 + z + \cdots$),
        then apply simple or higher-order pole formula as appropriate.
\end{enumerate}

\begin{learnbox}{Key Idea --- Section 4}
Always determine the pole order \emph{before} picking a formula.  A second-order pole treated
as first-order gives a wrong answer with no error message.
\end{learnbox}

% ================================================================
\section{Worked Examples}
% ================================================================

% ------------------------------------------------------------------
\subsection*{Example 1: Simple Pole at a Real Point}

Find $\operatorname{Res}\!\left(\dfrac{3z+1}{z-2},\; 2\right)$.

\textbf{Solution.}
$z = 2$ is a simple pole.  Using Formula 2:
\[
   \operatorname{Res} = \lim_{z\to 2}(z-2)\cdot\frac{3z+1}{z-2}
   = \lim_{z\to 2}(3z+1) = 3(2)+1 = 7.
\]

\begin{learnbox}{Takeaway --- Example 1}
For a simple pole of the form $\dfrac{g(z)}{z-z_0}$ with $g$ analytic at $z_0$, the residue
is simply $g(z_0)$.
\end{learnbox}

% ------------------------------------------------------------------
\subsection*{Example 2: Double Pole}

Find $\operatorname{Res}\!\left(\dfrac{e^z}{(z-1)^2},\; 1\right)$.

\textbf{Solution.}
$z = 1$ is a pole of order $m = 2$.  Using Formula 3:
\[
   \operatorname{Res} = \frac{1}{(2-1)!}\lim_{z\to 1}\frac{d}{dz}\!\left[(z-1)^2\cdot\frac{e^z}{(z-1)^2}\right]
   = \lim_{z\to 1}\frac{d}{dz}(e^z)
   = \lim_{z\to 1} e^z = e^1 = e.
\]

\begin{learnbox}{Takeaway --- Example 2}
For a double pole $\dfrac{g(z)}{(z-z_0)^2}$ with $g$ analytic at $z_0$, the residue equals $g'(z_0)$.
\end{learnbox}

% ------------------------------------------------------------------
\subsection*{Example 3: Rational Function with Factorised Denominator}

Find $\operatorname{Res}\!\left(\dfrac{z^2+2}{(z-1)(z+3)},\; 1\right)$ and
$\operatorname{Res}\!\left(\dfrac{z^2+2}{(z-1)(z+3)},\; -3\right)$.

\textbf{Solution.} Both poles are simple.

At $z = 1$ (Formula 4 with $q(z) = (z-1)(z+3)$, $q'(z) = 2z+2$):
\[
   \operatorname{Res}(f, 1) = \frac{1^2+2}{2(1)+2} = \frac{3}{4}.
\]

At $z = -3$ ($q'(z)$ evaluated at $-3$: $2(-3)+2 = -4$):
\[
   \operatorname{Res}(f,-3) = \frac{(-3)^2+2}{-4} = \frac{11}{-4} = -\frac{11}{4}.
\]

\begin{learnbox}{Takeaway --- Example 3}
When the denominator is already factorised, the $p/q'$ formula computes each simple-pole
residue in one arithmetic step.
\end{learnbox}

% ------------------------------------------------------------------
\subsection*{Example 4: Residue from Laurent Expansion Directly}

Find $\operatorname{Res}\!\left(\dfrac{\sin z}{z^4},\; 0\right)$.

\textbf{Solution.}
Expand $\sin z$ about $z = 0$:
\[
   \sin z = z - \frac{z^3}{3!} + \frac{z^5}{5!} - \cdots
\]
Dividing by $z^4$:
\[
   \frac{\sin z}{z^4} = z^{-3} - \frac{1}{6}\,z^{-1} + \frac{z}{120} - \cdots
\]
The coefficient of $z^{-1}$ is $a_{-1} = -\dfrac{1}{6}$.

Therefore $\operatorname{Res}\!\left(\dfrac{\sin z}{z^4},\, 0\right) = -\dfrac{1}{6}$.

\begin{learnbox}{Takeaway --- Example 4}
When the function involves a known Taylor series in the numerator, expanding and dividing is
often faster than applying the derivative formula.
\end{learnbox}

% ------------------------------------------------------------------
\subsection*{Example 5: Trigonometric Numerator --- Simple Pole}

Find $\operatorname{Res}\!\left(\dfrac{\cos z}{z - \pi/2},\; \pi/2\right)$.

\textbf{Solution.}
$z_0 = \pi/2$ is a simple pole.
\[
   \operatorname{Res} = \lim_{z\to\pi/2}(z - \pi/2)\cdot\frac{\cos z}{z - \pi/2}
   = \cos(\pi/2) = 0.
\]

\begin{learnbox}{Takeaway --- Example 5}
A simple pole with residue zero means the function is actually analytic at that point (the
apparent singularity is removable).  Always verify whether the apparent singularity is genuine.
\end{learnbox}

% ------------------------------------------------------------------
\subsection*{Example 6: Pole of Order 3}

Find $\operatorname{Res}\!\left(\dfrac{1}{z^3(z-2)},\; 0\right)$.

\textbf{Solution.}
$z = 0$ is a pole of order $m = 3$.
\[
   \phi(z) = z^3\cdot\frac{1}{z^3(z-2)} = \frac{1}{z-2}.
\]
\[
   \phi'(z) = -\frac{1}{(z-2)^2}, \qquad
   \phi''(z) = \frac{2}{(z-2)^3}.
\]
\[
   \operatorname{Res} = \frac{1}{2!}\lim_{z\to 0}\phi''(z)
   = \frac{1}{2}\cdot\frac{2}{(0-2)^3}
   = \frac{1}{2}\cdot\frac{2}{-8}
   = -\frac{1}{8}.
\]

\begin{learnbox}{Takeaway --- Example 6}
For a pole of order $m$, multiply out $(z-z_0)^m$, differentiate $(m-1)$ times, divide by $(m-1)!$,
and only then substitute $z = z_0$.  Do not rush to substitute earlier.
\end{learnbox}

% ------------------------------------------------------------------
\subsection*{Example 7: Exponential-Type Function}

Find $\operatorname{Res}\!\left(\dfrac{e^{2z}}{(z+1)^2},\; -1\right)$.

\textbf{Solution.}
$z = -1$ is a double pole ($m = 2$).
\[
   \phi(z) = (z+1)^2\cdot\frac{e^{2z}}{(z+1)^2} = e^{2z}.
\]
\[
   \frac{d}{dz}e^{2z} = 2e^{2z}.
\]
\[
   \operatorname{Res} = \frac{1}{1!}\lim_{z\to -1} 2e^{2z} = 2e^{-2}.
\]

\begin{learnbox}{Takeaway --- Example 7}
Exponential functions differentiate cleanly.  For double poles with exponential numerators,
the residue is just the first derivative of the numerator evaluated at the pole.
\end{learnbox}

% ------------------------------------------------------------------
\subsection*{Example 8: Rational Function --- $p/q'$ Method}

Find $\operatorname{Res}\!\left(\dfrac{z^3}{z^4-1},\; i\right)$.

\textbf{Solution.}
$z^4 - 1 = (z-1)(z+1)(z-i)(z+i)$.  So $z = i$ is a simple pole.
Using $p/q'$ with $p(z) = z^3$ and $q(z) = z^4 - 1$:
\[
   q'(z) = 4z^3, \qquad q'(i) = 4i^3 = 4(-i) = -4i.
\]
\[
   \operatorname{Res}(f,\, i) = \frac{i^3}{-4i} = \frac{-i}{-4i} = \frac{1}{4}.
\]

\begin{learnbox}{Takeaway --- Example 8}
For simple poles of the form $z^n/(z^{2n}-1)$ or similar symmetric rational functions, the
$p/q'$ formula gives the residue with minimal computation.
\end{learnbox}

% ------------------------------------------------------------------
\subsection*{Example 9: Mixed --- Simple Poles via Partial Fractions Check}

Find $\operatorname{Res}\!\left(\dfrac{2z+3}{z^2+3z+2},\; -1\right)$ and
$\operatorname{Res}\!\left(\dfrac{2z+3}{z^2+3z+2},\; -2\right)$.

\textbf{Solution.}
Factor: $z^2+3z+2 = (z+1)(z+2)$.  Both poles are simple.

$q'(z) = 2z+3$.

At $z = -1$:
\[
   \operatorname{Res}(f,-1) = \frac{2(-1)+3}{2(-1)+3} = \frac{1}{1} = 1.
\]

At $z = -2$:
\[
   \operatorname{Res}(f,-2) = \frac{2(-2)+3}{2(-2)+3} = \frac{-1}{-1} = 1.
\]

\textbf{Verification by partial fractions:}
\[
   \frac{2z+3}{(z+1)(z+2)} = \frac{A}{z+1} + \frac{B}{z+2}.
\]
Setting $z=-1$: $A = 1$.  Setting $z=-2$: $B = 1$.
The residues are indeed 1 and 1, confirming the $p/q'$ results. \checkmark

\begin{learnbox}{Takeaway --- Example 9}
Partial fractions provide an independent verification of residue computations for rational functions.
A residue equals the numerator coefficient in the partial fraction expansion.
\end{learnbox}

% ================================================================
\section{Method Selection Table (MANDATORY UPGRADE)}
% ================================================================

\begin{center}
\begin{tabular}{@{}>{\raggedright}p{3.2cm} >{\raggedright}p{3.5cm} >{\raggedright\arraybackslash}p{7cm}@{}}
\toprule
\textbf{Singularity Type} & \textbf{Best Method} & \textbf{Formula to Use} \\
\midrule
Simple pole, general & Direct limit &
$\operatorname{Res}(f,z_0) = \lim_{z\to z_0}(z-z_0)f(z)$ \\
\midrule
Simple pole, rational $p/q$ & $p/q'$ shortcut &
$\operatorname{Res}(f,z_0) = p(z_0)/q'(z_0)$ \\
\midrule
Pole of order $m \geq 2$ & Derivative formula &
$\dfrac{1}{(m-1)!}\lim_{z\to z_0}\dfrac{d^{m-1}}{dz^{m-1}}\!\left[(z-z_0)^m f(z)\right]$ \\
\midrule
Essential singularity & Laurent expansion &
Read $a_{-1}$ directly from the series \\
\midrule
Unclear order & Laurent expansion &
Expand until $(z-z_0)^{-1}$ term is identified; its coefficient is the residue \\
\midrule
Trig/exponential in numerator & Taylor expansion first &
Substitute series (e.g.\ $\sin z = z - z^3/6 + \cdots$), then divide and read $a_{-1}$ \\
\midrule
Removable singularity & None needed &
$\operatorname{Res}(f,z_0) = 0$ (no $z^{-1}$ term exists in Laurent series) \\
\bottomrule
\end{tabular}
\end{center}

\begin{learnbox}{Key Idea --- Section 6}
The method-selection table is a diagnostic checklist.  Step one is \emph{always} to determine
the type and order of the singularity; only then pick the matching formula.
\end{learnbox}

% ================================================================
\section{Excel Example (MANDATORY)}
% ================================================================

\textbf{Goal.} Use a spreadsheet to numerically estimate the residue of
$f(z) = \dfrac{1}{z(z-3)}$ at the simple pole $z = 3$ by computing
$(z-3)f(z) = \dfrac{1}{z}$ for values of $z$ approaching 3 from both sides.
The exact residue is $\operatorname{Res}(f,3) = \dfrac{1}{3}$.

\begin{center}
\begin{tabular}{@{}ccccc@{}}
\toprule
\textbf{Col A} & \textbf{Col B} & \textbf{Col C} & \textbf{Col D} & \textbf{Col E} \\
$z$ (real) & $z - 3$ & $f(z) = 1/(z(z-3))$ & $(z-3) \cdot f(z)$ & $|D - 1/3|$ (error) \\
\midrule
3.1  & \texttt{=A2-3}  & \texttt{=1/(A2*(A2-3))}  & \texttt{=B2*C2}  & \texttt{=ABS(D2-1/3)} \\
3.01 & \texttt{=A3-3}  & \texttt{=1/(A3*(A3-3))}  & \texttt{=B3*C3}  & \texttt{=ABS(D3-1/3)} \\
3.001& \texttt{=A4-3}  & \texttt{=1/(A4*(A4-3))}  & \texttt{=B4*C4}  & \texttt{=ABS(D4-1/3)} \\
\midrule
3.1  & 0.1    & 3.226... & 0.3226 & 0.0107 \\
3.01 & 0.01   & 33.22... & 0.3322 & 0.0011 \\
3.001& 0.001  & 333.2... & 0.3332 & 0.0001 \\
\bottomrule
\end{tabular}
\end{center}

\textbf{Spreadsheet instructions:}
\begin{enumerate}[leftmargin=*]
  \item Enter the $z$ values $3.1, 3.01, 3.001, 2.9, 2.99, 2.999$ in column A (rows 2--7).
  \item Column B: \texttt{=A2-3} (the ``approach distance'').
  \item Column C: \texttt{=1/(A2*(A2-3))} --- the full function $f(z)$.
  \item Column D: \texttt{=B2*C2} --- this cancels the simple pole and should converge to $1/3$.
  \item Column E: \texttt{=ABS(D2-1/3)} --- tracks the numerical error.
  \item In cell G2 type the analytical residue: \texttt{=1/3}.
  \item Observe that column D converges to $0.3333\ldots$ as $z \to 3$, confirming
        $\operatorname{Res}(f,3) = 1/3$.
\end{enumerate}

\begin{learnbox}{Key Idea --- Excel Section}
Multiplying $f(z)$ by $(z-z_0)$ and taking the numerical limit is the spreadsheet
analogue of the simple-pole residue formula.  Column D converging to a constant is
strong evidence of a simple pole with that constant as its residue.
\end{learnbox}

% ================================================================
\section{Python Example (MANDATORY)}
% ================================================================

The following script uses \texttt{sympy} to compute residues for several
representative functions from this topic.

\begin{lstlisting}[language=Python]
from sympy import *

z = symbols('z')

# Define functions
f1 = (3*z + 1) / (z - 2)
f2 = exp(z) / (z - 1)**2
f3 = (z**2 + 2) / ((z - 1)*(z + 3))
f4 = sin(z) / z**4
f5 = (z**3) / (z**4 - 1)
f6 = 1 / (z**3 * (z - 2))
f7 = exp(2*z) / (z + 1)**2
f8 = (2*z + 3) / (z**2 + 3*z + 2)

functions = [
    (f1, 2,   "Res( (3z+1)/(z-2),      z=2  )"),
    (f2, 1,   "Res( exp(z)/(z-1)^2,    z=1  )"),
    (f3, 1,   "Res( (z^2+2)/(...),      z=1  )"),
    (f3, -3,  "Res( (z^2+2)/(...),      z=-3 )"),
    (f4, 0,   "Res( sin(z)/z^4,         z=0  )"),
    (f5, I,   "Res( z^3/(z^4-1),        z=i  )"),
    (f6, 0,   "Res( 1/(z^3(z-2)),       z=0  )"),
    (f7, -1,  "Res( exp(2z)/(z+1)^2,   z=-1 )"),
    (f8, -1,  "Res( (2z+3)/(z^2+3z+2), z=-1 )"),
    (f8, -2,  "Res( (2z+3)/(z^2+3z+2), z=-2 )"),
]

print("=" * 55)
print("RESIDUE COMPUTATIONS USING SymPy")
print("=" * 55)
for (func, point, label) in functions:
    res = residue(func, z, point)
    print(f"{label}")
    print(f"  = {res}")
    print()

# Expected output:
# Res( (3z+1)/(z-2),      z=2  )  =  7
# Res( exp(z)/(z-1)^2,    z=1  )  =  E  (i.e., e)
# Res( (z^2+2)/(...),      z=1  )  =  3/4
# Res( (z^2+2)/(...),      z=-3 )  =  -11/4
# Res( sin(z)/z^4,         z=0  )  =  -1/6
# Res( z^3/(z^4-1),        z=i  )  =  1/4
# Res( 1/(z^3(z-2)),       z=0  )  =  -1/8
# Res( exp(2z)/(z+1)^2,   z=-1 )  =  2*exp(-2)
# Res( (2z+3)/(z^2+3z+2), z=-1 )  =  1
# Res( (2z+3)/(z^2+3z+2), z=-2 )  =  1

# Visualise the magnitude of f1 near z=2 on the real axis
import numpy as np
import matplotlib
matplotlib.use('Agg')
import matplotlib.pyplot as plt

x_vals = np.linspace(1.5, 2.5, 400)
# Avoid the pole point exactly
mask = np.abs(x_vals - 2) > 1e-3
f1_vals = (3*x_vals[mask] + 1) / (x_vals[mask] - 2)
plt.figure()
plt.plot(x_vals[mask], f1_vals, color='blue')
plt.axvline(x=2, color='red', linestyle='--', label='pole z=2')
plt.xlabel('x (real axis)')
plt.ylabel('f1(x) = (3x+1)/(x-2)')
plt.title('Behaviour of f1 near simple pole z=2')
plt.legend()
plt.ylim(-60, 60)
plt.grid(True)
plt.savefig('residue_pole_plot.png', dpi=100, bbox_inches='tight')
print("Plot saved to residue_pole_plot.png")
\end{lstlisting}

\begin{learnbox}{Key Idea --- Python Section}
\texttt{sympy.residue(f, z, z0)} computes the exact symbolic residue.
For simple poles it applies the limit formula; for higher-order poles it uses the derivative
formula automatically.  Always verify SymPy results against hand calculations on at least
one example.
\end{learnbox}

% ================================================================
\section{Viva Questions (8)}
% ================================================================

\begin{enumerate}[leftmargin=*]
  \item \textbf{What is the residue of $f$ at an isolated singularity $z_0$?}\\
        It is the coefficient $a_{-1}$ in the Laurent expansion of $f$ about $z_0$,
        i.e.\ $\operatorname{Res}(f,z_0) = a_{-1}$.

  \item \textbf{State the simple-pole residue formula.}\\
        If $z_0$ is a simple pole, $\operatorname{Res}(f,z_0) = \lim_{z\to z_0}(z-z_0)f(z)$.

  \item \textbf{How does the residue formula change for a pole of order $m$?}\\
        $\operatorname{Res}(f,z_0) = \dfrac{1}{(m-1)!}\lim_{z\to z_0}
        \dfrac{d^{m-1}}{dz^{m-1}}\!\left[(z-z_0)^m f(z)\right].$

  \item \textbf{When is the formula $\operatorname{Res}(p/q,\, z_0) = p(z_0)/q'(z_0)$ valid?}\\
        When $q(z_0)=0$, $q'(z_0)\neq 0$, and $p(z_0)\neq 0$ --- i.e.\ $z_0$ is a simple
        zero of $q$ and not a zero of $p$.

  \item \textbf{What is the residue at a removable singularity?}\\
        Zero, because there is no $z^{-1}$ term in the Laurent expansion at a removable singularity.

  \item \textbf{Why is the residue independent of which Laurent series annulus is used?}\\
        All Laurent series of $f$ in annuli sharing the same inner singularity $z_0$ have the same
        $a_{-1}$ coefficient, since the contour-integral formula for $a_{-1}$ depends only on $f$
        and $z_0$, not on the outer radius.

  \item \textbf{Can a function have the same residue at two different poles?}\\
        Yes.  For example, $f(z) = \dfrac{2z+3}{(z+1)(z+2)}$ has residue 1 at both $z=-1$ and $z=-2$.

  \item \textbf{What is the connection between residues and contour integrals?}\\
        By the Residue Theorem, $\oint_C f(z)\,dz = 2\pi i \sum \operatorname{Res}(f, z_k)$
        where the sum is over all isolated singularities $z_k$ enclosed by $C$.
\end{enumerate}

% ================================================================
\section{Descriptive Questions (5)}
% ================================================================

\begin{enumerate}[leftmargin=*]
  \item Define the residue of a complex function at an isolated singularity.  Derive the residue
        formula for a pole of order $m$ starting from the Laurent series and the contour-integral
        formula for Laurent coefficients.

  \item State and derive the formula $\operatorname{Res}(p/q, z_0) = p(z_0)/q'(z_0)$ for a simple
        pole of $p(z)/q(z)$.  Use L'H\^{o}pital's rule or direct factorisation in your derivation.

  \item Find the residues of $f(z) = \dfrac{z^2 - 4z + 5}{(z-1)^2(z-3)}$ at all its singularities.
        Show every step and specify the order of each pole before applying any formula.

  \item Explain, with an example, why Laurent series expansion is the most general method for computing
        residues.  Include an example involving an essential singularity (e.g.\ $e^{1/z}$ at $z=0$).

  \item Compute $\operatorname{Res}\!\left(\dfrac{e^{3z}}{(z+2)^3},\; -2\right)$ using the
        order-3 pole formula.  Verify your answer by expanding $e^{3z}$ in a Taylor series
        about $z=-2$ and reading off the $a_{-1}$ coefficient from $f$'s Laurent expansion.
\end{enumerate}

% ================================================================
\section{Practice Problems (6)}
% ================================================================

\begin{enumerate}[leftmargin=*]
  \item Find $\operatorname{Res}\!\left(\dfrac{5z-2}{z(z-1)},\; 0\right)$ and
        $\operatorname{Res}\!\left(\dfrac{5z-2}{z(z-1)},\; 1\right)$.
        \textit{[Hint: simple poles; use $p/q'$.  Answers: $-(-2)/(-1) = 2$ at $z=0$;
        $(5-2)/1 = 3$ at $z=1$.]}

  \item Find $\operatorname{Res}\!\left(\dfrac{\cos z}{z^2},\; 0\right)$.
        \textit{[Hint: expand $\cos z = 1 - z^2/2 + \cdots$, divide by $z^2$; $a_{-1} = 0$.]}

  \item Find $\operatorname{Res}\!\left(\dfrac{z}{(z^2+4)^2},\; 2i\right)$.
        \textit{[Hint: double pole at $2i$; factor $(z^2+4)^2 = (z-2i)^2(z+2i)^2$; answer $= -i/16$.]}

  \item Find $\operatorname{Res}\!\left(\dfrac{e^z}{z^4},\; 0\right)$.
        \textit{[Hint: $e^z = 1 + z + z^2/2 + z^3/6 + \cdots$; divide by $z^4$;
        $a_{-1}$ is the $z^{-1}$ term's coefficient $= 1/6$.]}

  \item Find $\operatorname{Res}\!\left(\dfrac{1}{z^2(z^2+1)},\; 0\right)$.
        \textit{[Hint: expand $1/(z^2+1) = 1 - z^2 + z^4 - \cdots$ for $|z|<1$;
        then $f = z^{-2}(1 - z^2 + \cdots) = z^{-2} - 1 + \cdots$; $a_{-1} = 0$.]}

  \item Find all residues of $f(z) = \dfrac{z^2}{(z-1)(z-2)(z-3)}$.
        \textit{[Hint: three simple poles; use $p/q'$.  At $z=k$: $q'(k)$ is the product of
        $(k-j)$ for $j\neq k$.  Answers: $1/2$ at $z=1$; $-4$ at $z=2$; $9/2$ at $z=3$.]}
\end{enumerate}

% ================================================================
\section{MCQs (5)}
% ================================================================

\begin{enumerate}[leftmargin=*]

\item The residue of $f(z) = \dfrac{1}{z-4}$ at $z=4$ is:
\begin{enumerate}[(a)]
  \item $0$
  \item $4$
  \item \textbf{$1$}
  \item $-1$
\end{enumerate}
\textit{Explanation:} $z=4$ is a simple pole; $\operatorname{Res} = \lim_{z\to 4}(z-4)/(z-4) = 1$.

\item The residue of $\dfrac{e^z}{z^2}$ at $z=0$ is:
\begin{enumerate}[(a)]
  \item $0$
  \item \textbf{$1$}
  \item $e$
  \item $2$
\end{enumerate}
\textit{Explanation:} $e^z/z^2 = z^{-2} + z^{-1} + 1/2 + \cdots$; the coefficient of $z^{-1}$ is $1$.

\item If $z_0$ is a pole of order $m=3$ of $f$, the residue formula requires differentiating
$(z-z_0)^3 f(z)$ how many times?
\begin{enumerate}[(a)]
  \item 3 times
  \item \textbf{2 times}
  \item 1 time
  \item 0 times
\end{enumerate}
\textit{Explanation:} The formula uses the $(m-1) = 2$nd derivative.

\item The residue of $\dfrac{\sin z}{z^3}$ at $z=0$ is:
\begin{enumerate}[(a)]
  \item $1$
  \item $0$
  \item \textbf{$-1/6$}
  \item $1/6$
\end{enumerate}
\textit{Explanation:} $\sin z/z^3 = z^{-2} - (1/6)z^{-0} + \cdots$; wait --- $\sin z = z - z^3/6 + \cdots$, so $\sin z / z^3 = z^{-2} - 1/6 + z^2/120 - \cdots$; there is no $z^{-1}$ term, so $\operatorname{Res} = 0$.

\textit{Correction:} Residue $= 0$.  (See option (b).)

\textbf{Corrected answer: (b) $0$.}

\textit{Explanation:} $\sin z = z - z^3/6 + z^5/120 - \cdots$; dividing by $z^3$ gives
$z^{-2} - 1/6 + z^2/120 - \cdots$.  No $z^{-1}$ term; residue $= 0$.

\item For $f(z) = p(z)/q(z)$ with $q(z_0)=0$, $q'(z_0)\neq 0$, the residue at $z_0$ equals:
\begin{enumerate}[(a)]
  \item $p'(z_0)/q(z_0)$
  \item $p(z_0)/q(z_0)$
  \item \textbf{$p(z_0)/q'(z_0)$}
  \item $p'(z_0)/q'(z_0)$
\end{enumerate}
\textit{Explanation:} This is the standard $p/q'$ formula for a simple pole of a rational function.

\end{enumerate}

% ================================================================
\section{Common Mistakes}
% ================================================================

\begin{mistakebox}{Common Mistakes When Computing Residues}
\begin{tabular}{@{}>{\raggedright}p{3.8cm} >{\raggedright}p{4.5cm} >{\raggedright\arraybackslash}p{5.5cm}@{}}
\toprule
\textbf{Mistake} & \textbf{What Goes Wrong} & \textbf{Correct Approach} \\
\midrule
Not identifying pole order first & Applying the simple-pole formula to a double pole gives a wrong answer & Always factor the denominator or check the Laurent series to determine $m$ before choosing a formula \\
\midrule
Sign error after partial fractions & Forgetting the negative sign when $q'(z_0) < 0$ & Carefully substitute $z_0$ into $q'$ and carry the sign through the fraction \\
\midrule
Wrong derivative order in higher-order pole formula & Differentiating $m$ times instead of $m-1$ times & The formula requires the $(m-1)$th derivative of $(z-z_0)^m f(z)$ \\
\midrule
Omitting division by $(m-1)!$ & The formula includes a $1/(m-1)!$ factor that students forget & Write the full formula first, then evaluate: $\dfrac{1}{(m-1)!}\cdot\lim(\cdots)$ \\
\midrule
Substituting $z=z_0$ before differentiating & Produces $0/0$ or an undefined expression & Differentiate $(z-z_0)^m f(z)$ symbolically, then substitute $z_0$ \\
\midrule
Treating a removable singularity as a pole & Applying residue formulas where the function is actually analytic & If $(z-z_0)f(z)\to 0$ as $z\to z_0$, the singularity is removable; residue $=0$ \\
\midrule
Using $\int$ instead of $\oint$ for closed contours & Notation inconsistency; may confuse orientation conventions & Always write $\oint_C$ for closed contour integrals \\
\bottomrule
\end{tabular}
\end{mistakebox}

% ================================================================
\section{Quick Recap}
% ================================================================

\begin{learnbox}{Quick Recap --- Topic 17: Residues}
\begin{itemize}[leftmargin=*]
  \item The residue $\operatorname{Res}(f,z_0) = a_{-1}$ is the coefficient of $(z-z_0)^{-1}$ in the Laurent series of $f$ about $z_0$.
  \item For a simple pole: $\operatorname{Res}(f,z_0) = \lim_{z\to z_0}(z-z_0)f(z)$.
  \item For a rational function $p/q$ with a simple zero of $q$ at $z_0$: $\operatorname{Res} = p(z_0)/q'(z_0)$.
  \item For a pole of order $m$: $\operatorname{Res}(f,z_0) = \dfrac{1}{(m-1)!}\lim_{z\to z_0}\dfrac{d^{m-1}}{dz^{m-1}}\!\left[(z-z_0)^m f(z)\right]$.
  \item For functions with essential singularities or unclear order: expand in Laurent series and read $a_{-1}$ directly.
  \item The residue at a removable singularity is always zero.
  \item Always determine the pole order \emph{before} choosing a residue formula.
  \item Residues are the key ingredient in the Cauchy Residue Theorem (Topic 18), enabling closed-contour integrals to be evaluated by pure algebra.
\end{itemize}
\end{learnbox}

\end{document}
